\section{预期结果与验收标准}
\label{sec:results}
\subsection{匿名化验收}
预期匿名结果只改变面部内部，保留头发、头部外轮廓及座舱背景。验收拟同时查看局部对照和自建身份模型的预测变化，目标是在表情结构可用的条件下压低原身份命中。掩膜外灰度像素保持作为单独检查项，确认处理范围符合设计。

\subsection{表情识别验收}
表情模型输入必须来自匿名结果，以真实标签统计准确率与各类召回。拟分别评价通用七类表情、车载六类表情及近红外中性样本，并观察中性与悲伤、愤怒之间的混淆。候选比较保留完整集合与固定配置，使优化前后可以直接核对。

\subsection{实时性能验收}
性能验收统计模型完成匿名化和分类所需的时间，分别观察启动和稳态表现。连续视频中检查切换表情、停止处理、来源切换与失脸恢复，确认当前结果对应当前画面。质量配置优先保证准确性，高速配置再以接近视频连续观察需求为优化目标。

\subsection{成果形式}
拟交付三页桌面应用、可复现源代码、必要模型资源说明、项目计划书、研究报告及汇报材料。演示中用带标签图片说明匿名化与表情保持，用示例视频和摄像头说明连续分析，用分域图表说明整体表现。
